\chapter{Forward and Futures Markets}

These notes follow Robert J. Shiller's lecture on forward and futures markets, curated by LazyingArt LLC for the Yale lecture series. The lecture begins by enlarging the stakes. Futures markets matter, he says, because they give us prices for the future, and serious planning requires a long horizon. He even puts the point in personal terms: if we plan only for the present, we are planning on too small a scale. From that opening the lecture moves, in order, through forwards and derivatives, public hostility to speculation, grain storage, the historical invention of futures trading at Dojima, the modern agricultural term structure, margin and daily settlement, the fair-value storage relation, oil backwardation, convenience yield, and finally stock-index futures. The whole chapter is guided by one idea: a futures market is not merely a wager on a later date. It is an organized way of pricing time, storage, and risk.

\section{Prices for the future}

There are two ways to look at a futures market, and the lecture deliberately begins with the broader one. Technically, a futures contract is a standardized contract traded in an organized market for delivery at a future date. More broadly, futures markets are markets about the future. They are places where guesses about later conditions are translated into current prices.

That is why the lecture links the subject back to earlier material on forwards and forward rates. Forwards and futures are related ideas, but futures are the more developed institutional form. A futures market resembles a stock exchange: the contract is standardized, widely observed, and publicly quoted. A forward market is different. It is a bilateral arrangement between particular counterparties, with custom terms about quality, location, timing, and failure. Both are derivatives in the lecture's sense, because their price derives from an underlying market.

The early motivational detour into derivatives is not incidental. Shiller pauses over the public hostility to the word because he wants to clear away a conceptual obstacle before the agricultural examples begin. To people inside finance, derivative markets are an extension of information: they take an underlying market and add more detail to it. To the broader public, however, the word often carries suspicion and anger. The lecture says that regulation is necessary, both private and governmental, but it also insists that there is nothing inherently evil about derivative markets. In a modern economy they are part of how we do the arithmetic of the future.

This is why he invokes Charles Conant's old point about speculation. Business decisions always involve guesses about what lies ahead. In a well-developed market, those guesses become prices, and those prices then enter the calculations made by households, firms, and governments. A futures market is therefore not an optional extra laid on top of ``real'' economic life. It is one of the institutions by which real economic life thinks ahead.

That claim is easiest to see in agriculture, where time and storage are impossible to ignore.

\section{Speculation, storage, and the grain market}

The lecture begins its substantive analysis with a question from a survey once given in New York and Moscow. If grain traders hold grain in storage in anticipation of higher prices later, do they make shortages more common, or rarer? The answer given by much of the public was the wrong one. In the lecture's retelling, many respondents treated speculation in grain as something that creates shortages. That, he argues, reverses the logic of storage.

In most places there is one main wheat harvest each year. People do not consume the whole crop immediately. So someone must carry grain from harvest season to later months. Warehouses and grain merchants are not performing a morally dubious side activity; they are solving a basic intertemporal problem. Somewhere, as the lecture puts it, there is a warehouse storing the grain that will be consumed six months from now.

Now suppose those professionals think grain will be scarcer later. They hold more back now. The present price rises. That higher present price reduces present consumption and leaves more inventory for the future. In this sense expected scarcity tomorrow can justify higher prices today, and the higher current price is precisely what smooths use over time.

\subsection{Question \& Answer}

\paragraph{Question.} Does speculation in grain make shortages worse, or can it actually stabilize supply?

\paragraph{Answer.} In the lecture's logic it can stabilize supply. If expected future scarcity induces storage now, the current price rises. That higher price causes current users to consume a little less, and more grain remains available for later delivery. Storage in anticipation of higher later prices is therefore not necessarily an antisocial act. In a seasonal commodity market it is one of the mechanisms by which scarcity is spread more evenly through time.

This is the lecture's first genuine mechanism, and it is worth stating plainly: storage decisions today affect prices today, consumption today, and shortages tomorrow. Futures markets do not create that mechanism. They sharpen it.

The lecture then turns this economic point into a moral one. When we see a strong incentive to store grain, we may be watching people act out of self-interest, but the effect can still be socially useful. This is the background for the later argument that a steep contango in grain is not merely a prediction of future prices. It is an inducement to expand storage.

\section{From Dojima to the centralized exchange}

Once we understand why storage matters, the historical move to Dojima becomes natural. The lecture takes us to Osaka in 1673, where rice was central to economic life and Dojima served as a national rice market. Shiller cites a figure of $91$ rice warehouses in Dojima in that year. The commodity, the storage system, and the need for future delivery were already in place.

Before a true futures exchange existed, merchants used forward contracts. A rice merchant in some nearby town could arrange with a warehouse for regular delivery on agreed terms. Forward contracting is therefore the first form of future delivery, and the lecture says so explicitly: forwards precede futures.

But the forward market developed problems that called for institutional repair. The first problem was counterparty risk. If the market moved against one side, that side might simply fail to perform. If rice prices fell, the buyer might ignore the contract and purchase more cheaply elsewhere. If prices rose, the warehouse might try to escape an unfavorable deal. The second problem was heterogeneity. Every contract was different: different quality, different place of delivery, different remedies, different conditions. If someone announced a ``price of rice,'' one had to ask immediately: which rice, where, under what terms? The market produced promises, but not yet a clean public price.

Dojima's innovation was to create standardized exchange trading. Contracts were mediated by the exchange and traded during specified hours on a common floor. The details the lecture emphasizes are telling: a fuse was reportedly burned down to mark the end of trading; ``water men'' enforced the close; hand signals arose because floor trading was too loud for speech alone. These are not merely picturesque details. They show us a market becoming an institution.

In modern language, the same principle survives. A futures contract specifies the commodity, the delivery date, the approved warehouse, and the inspection regime. Inspectors determine quality according to exchange standards. No single trader gets precisely customized terms, but everyone gains a contract that is legible, comparable, and tradable.

\subsection{Question \& Answer}

\paragraph{Question.} What exact problem did futures markets solve that forward contracts could not?

\paragraph{Answer.} They solved several linked problems at once. They replaced a mass of idiosyncratic bilateral contracts with a standardized object, thereby making price quotations meaningful. They sharply reduced counterparty risk by routing performance through the exchange structure rather than leaving it to the honor or solvency of a specific counterparty. And they created a centralized market in which contracts could be traded rather than merely held until delivery. The lecture's point is that futures trading did not arise because traders wanted a more abstract market. It arose because the forward market was too opaque and too fragile.

This historical point leads directly into the modern one. The lecture notes that the inventors of rice futures no longer host the main rice futures market. Modern rice futures are quoted in the United States, because a futures market must be centralized where the trading and standardization succeed. That centralized futures price then becomes, for many purposes, the market price that farmers, merchants, and processors actually use.

\section{Modern agricultural futures, contango, and daily settlement}

At this point the lecture returns from history to present-day agricultural markets. The move matters because Shiller wants us to see that a futures market often becomes the practical quoted market. Farmers hear wheat, soybean, or rice prices on the radio, and those prices are usually futures prices. The reason is not mystical. The futures contract is well defined, standardized, and centrally quoted, whereas the spot market is local, heterogeneous, and often hard to interpret.

Rough rice provides the first modern example. The lecture describes a futures curve with different delivery months extending roughly a year into the future, quoted in cents per hundredweight. The front end of the curve is the reliable part because that is where trading is concentrated. Farther-out contracts are thinner and less informative. The striking fact, however, is the slope: the curve rises from roughly $13.75$ to $15.5$ over about six months.

The first quantitative reaction is the right one:
\begin{equation}
\frac{15.5-13.75}{13.75}\approx 0.127.
\end{equation}
That is roughly a $12.7\%$ increase over half a year. If this were a free carry trade, it would be an obvious profit opportunity.

Now comes an important warning from the lecture. This curve is a cross section of contracts quoted today at different maturities. It is not a graph of a single contract moving through time. An upward-sloping cross section is what the lecture calls contango; a downward-sloping one is backwardation. These are deliberately simple definitions, and the lecturer says there are subtler ones, but for the present chapter the simple curve-shape meaning is the one that matters.

\subsection{Question \& Answer}

\paragraph{Question.} If the futures curve is steeply upward sloping, why is there not an obvious riskless profit from buying spot and selling futures?

\paragraph{Answer.} Because the trade is only apparently free. To exploit the spread one must buy the commodity, store it, inspect it, insure it, finance it, and carry it until delivery. The lecture insists that professionals in these markets already know exactly what the curve says. A strong contango therefore generates efforts to expand storage, not unclaimed riches lying on the table. It is a market signal telling warehouse operators to search for more storage capacity, more sanitation clearance, more insurance, and more inventory space.

This is why the lecture returns to the larger social claim. A strong contango in grain encourages exactly the activity that makes famine less likely: additional storage by competent, profit-seeking operators. The lecture explicitly connects this to Adam Smith's theme that self-interested market participants may contribute more to social stability than more visibly benevolent actors who do not understand the storage problem at all.

The lecture then pivots back to institutional design. Forward contracts are burdened by counterparty risk, but futures markets largely eliminate that burden by means of margin and daily settlement. A trader can post on the order of \$5{,}000, face a margin requirement of roughly $3\%$ to $5\%$, and thereby take on something like \$100{,}000 of notional futures exposure. The contract price is not paid in full today. Instead, the position is marked to market each day.

For a short position of quantity $Q$, the lecture's verbal description can be rendered as the note-writing reconstruction
\begin{equation}
M_{t+1}=M_t-Q\bigl(F_{t+1}-F_t\bigr),
\end{equation}
where $M_t$ is the trader's margin balance and $F_t$ is the futures price on day $t$. If the futures price rises, the short's margin account is debited. If it falls, the account is credited. If the margin is nearly depleted, the broker demands more collateral; if none is posted, the broker closes the position.

\subsection{Question \& Answer}

\paragraph{Question.} Why does daily settlement make a futures contract safer than a private forward contract?

\paragraph{Answer.} Because losses are realized continuously rather than stored up until the delivery date. The exchange system does not wait to discover at maturity whether one side can pay. It checks every day. If the account falls too far, the position is closed. In that sense each trader contracts with the exchange mechanism rather than relying on the personal reliability of an unknown counterparty. That is why the lecture says futures prices are ``pure'' prices: they are not heavily contaminated by bilateral default risk.

The wheat example, which follows, confirms the same logic and also prepares the mathematical turn. The front month of soft red winter wheat is around \$7.23 per bushel. Again we see sharp contango at short maturities, followed by flattening or even decline around the next harvest. The natural lecture interpretation is seasonal. Storage compensation makes sense before the new crop arrives, but one should not expect the same storage premium to carry cleanly through harvest.

At this point the lecture says, in effect: now that the puzzle is alive, let us write down the simple mathematics.

\section{Fair value, storage cost, and convergence}

The board equation appears only after the institutional and motivational groundwork has been laid. That order matters. Shiller does not begin with algebra. He begins with the storage problem, with Dojima, with the rice curve, and with the question of why contango is not free money. Only then does he write the carrying-cost relation.

\begin{figure}[h]
\centering
\includegraphics[width=0.78\textwidth]{lecture_15_figure_02.png}
\caption{Blackboard fair-value equation for futures pricing.}
\end{figure}

The board gives the basic fair-value relation as
\begin{equation}
P_F=(1+r+s)P_{\text{spot}}.
\end{equation}
The notation should be read carefully. $P_F$ is the futures price for a specified maturity. $P_{\text{spot}}$ is the spot price today. The term $r$ is the total interest cost between now and maturity, not an annualized rate quoted independently of horizon. The term $s$ is the storage cost over the same horizon.

The lecture is explicit about the meaning of $r$. If the relevant interest rate is $5\%$ per year, then for a one-year contract we have
\begin{equation}
r=0.05,
\end{equation}
whereas for a one-and-a-half-year contract we would have, in the lecture's simple arithmetic,
\begin{equation}
r\approx 0.075.
\end{equation}
The point is not the exact compounding convention. The point is that both $r$ and $s$ are horizon terms. The farther out the contract, the larger the carrying burden.

There is a small conceptual subtlety here which the lecture itself notes. In practice the spot price of a commodity may be difficult to observe cleanly, precisely because the futures market has become the best quoted market. So $P_{\text{spot}}$ has a slightly theoretical flavor. But the carrying-cost logic still survives as the simplest way to think about how a storable commodity should price across maturities.

\paragraph{Worked derivation.} The lecture's logic can be organized as follows. Buy the commodity today, finance the purchase, store the commodity, and sell a futures contract against the inventory. Then
\begin{align}
\text{purchase cost today} &= P_{\text{spot}},\\
\text{financing cost to maturity} &= rP_{\text{spot}},\\
\text{storage cost to maturity} &= sP_{\text{spot}},\\
\text{carried value at delivery} &= (1+r+s)P_{\text{spot}}.
\end{align}
If the futures price were materially above this level, storage plus a short futures position would offer an abnormal profit. If it were materially below it, ordinary storage would be unattractive. That is why the board labels the relation ``fair value.'' It is the lecture's name for the carrying-cost benchmark.

But the lecture also immediately weakens any temptation to treat the equation as exact. If it held with mechanical exactness, there would be no normal storage profit at all. So this is not a theorem that must hold point by point to the last cent. It is the benchmark around which a storage market organizes itself.

Shiller then restates the point in words: the normal case is contango, because looking farther into the future usually means paying more total interest and more total storage cost. The equation is, as he says, an expression of ``the value of the future.''

\subsection{Question \& Answer}

\paragraph{Question.} How can the futures price be above spot now and still converge to spot at expiration?

\paragraph{Answer.} Because the extra terms above spot are the costs of carrying the commodity through the remaining life of the contract. As expiration approaches, that remaining life shrinks. There is less interest still to be paid and less storage still to be incurred. Hence the wedge between futures and spot shrinks with time.

In a cleaned notation, the lecture's convergence claim can be written as
\begin{equation}
T-t\downarrow 0
\quad\Rightarrow\quad
r(T-t),\,s(T-t)\downarrow 0
\quad\Rightarrow\quad
P_F(t,T)\to P_{\text{spot}}(T).
\end{equation}

The lecture then slows down and redraws the idea in calendar time rather than maturity time. This is one of the most important didactic moves in the chapter, because it prevents us from confusing an upward-sloping curve across maturities with an upward time path of a single contract.

\begin{figure}[h]
\centering
\includegraphics[width=0.92\textwidth]{lecture_15_figure_03.png}
\caption{Fair-value formula beside the seasonal storage sketch.}
\end{figure}

The board sketch shows a yearly cycle. Within the year the spot price of wheat rises as inventories are carried. At harvest the price drops sharply when new supply arrives. Then the cycle repeats. In the perfect world imagined in the lecture, there is no harvest uncertainty at all. If we fix one particular futures contract whose maturity lies just before harvest, then its price is constant through calendar time: it always equals the spot price that will prevail on that maturity date. The moving object is the spot price, which rises toward that fixed maturity value as the remaining carrying cost is used up.

\begin{figure}[h]
\centering
\begin{tikzpicture}[scale=0.95]
\draw[->] (0,0) -- (11.1,0) node[right] {calendar time};
\draw[->] (0,0) -- (0,4.1) node[above] {$P$};

\draw (2.5,0) -- (2.5,3.1);
\draw (5.0,0) -- (5.0,3.1);
\draw (7.5,0) -- (7.5,3.1);

\node[below] at (2.5,0) {$1$};
\node[below] at (5.0,0) {$2$};
\node[below] at (7.5,0) {$3$};

\draw[thick] (0.15,1.0) .. controls (0.9,1.15) and (1.7,2.0) .. (2.45,3.0);
\draw[thick] (2.55,1.0) .. controls (3.3,1.15) and (4.2,2.0) .. (4.95,3.0);
\draw[thick] (5.05,1.0) .. controls (5.8,1.15) and (6.7,2.0) .. (7.45,3.0);
\draw[thick] (7.55,1.0) .. controls (8.3,1.15) and (9.2,2.0) .. (9.95,3.0);

\draw[dashed] (0.15,3.0) -- (9.95,3.0);
\node[right] at (10.0,3.0) {fixed-maturity futures price in the no-uncertainty benchmark};

\node at (1.2,0.45) {harvest};
\node at (3.75,0.45) {harvest};
\node at (6.25,0.45) {harvest};

\node[left] at (0.0,2.2) {spot};
\end{tikzpicture}
\caption{Clean reconstruction of the lecture's seasonal calendar-time sketch. The spot price rises through the storage season and resets downward at harvest; a fixed-maturity futures price is flat in the no-uncertainty benchmark.}
\end{figure}

The lecture is therefore making two claims at once. Across maturities at a single date we usually see contango because carrying costs are positive. Through calendar time, however, a fixed-maturity futures price converges to spot as the remaining carry shrinks to zero.

\section{Oil, backwardation, and convenience yield}

The lecture now performs its sharpest turn. After making the fair-value relation sound natural and persuasive, it places next to it the crude-oil futures curve and asks: how can both be true at once?

\begin{figure}[h]
\centering
\includegraphics[width=0.96\textwidth]{lecture_15_figure_04.png}
\caption{Board pricing relation next to the crude-oil futures curve.}
\end{figure}

The curve shown on the slide is not monotone. It rises into late 2011, then declines, bottoming near January 2015, and then rises again farther out. The front month is around \$101 per barrel. The market is therefore pricing an increase, then a decrease, then a recovery. Shiller explicitly reads this as a market absorbed in large events: Middle Eastern disruption, energy policy after Japan's earthquake and nuclear accident, and the broader world dependence on oil. He also notes that oil and food are linked, because grain can be turned into ethanol and thereby partially substitute for oil.

The crucial interpretive point is that the front-month futures price is, in practice, what people mean when they speak of ``the price of oil.'' The true spot market for oil is too heterogeneous and too thin to provide a single clean price. Oil is sold under many long-term contracts with numerous contingencies. The quoted price is therefore the nearest futures contract, precisely because that is the standardized and legible market.

Now the puzzle becomes acute. The front month is almost spot because it is close to expiration. Yet farther-out prices can lie below it. How can that happen if the board formula says financing and storage cost should push futures above spot?

\subsection{Question \& Answer}

\paragraph{Question.} How can a futures curve slope downward if the fair-value formula adds interest and storage cost to spot?

\paragraph{Answer.} Because the simple carrying-cost relation is most natural when the commodity is in storage as part of an ordinary carry trade. A downward-sloping curve tells us that, over that stretch of maturities, nobody plans to store the commodity purely for speculative carry. If one tried to store oil from now to 2015 as a normal storage business while the curve sloped downward, one would lose money. So the simple storage relation is no longer a complete description of the term structure.

The lecture first puts this somewhat heuristically: one might say that storage costs have become ``negative.'' But the more careful refinement comes with the term convenience yield. A factory that depends on oil does not want its tanks to run dry. Even if storing oil is not profitable as a carry trade, holding inventory still provides operational protection, a buffer against disruption, a service flow from inventory itself.

That logic is supported by the later board frame:

\begin{figure}[h]
\centering
\includegraphics[width=0.92\textwidth]{lecture_15_figure_05.png}
\caption{Fair-value equation revisited with convenience-yield discussion.}
\end{figure}

The board note for convenience yield is faint, so the final formula should be stated cautiously. A lecture-aligned reconstruction is
\begin{equation}
P_F \approx (1+r+s-y)P_{\text{spot}},
\end{equation}
where $y$ is the convenience yield. This is not a verbatim transcription of the board, but it is the natural mathematical completion of the lecture's verbal point. When $y$ is large enough, it can offset financing and storage costs and permit backwardation.

This yields the lecture's main caveat. The fair-value storage relation is not false, but it is conditional. It works most cleanly when the commodity is in storage as a carry asset. Once inventory has an independent operational value, or once storage incentives largely disappear, the whole futures curve need not be governed by the simple formula alone.

\section{Oil history and financial futures}

Having resolved the local oil puzzle, the lecture zooms out historically. This matters for the chapter's rhythm. The oil example is not left as a technical footnote. It is enlarged into a story about the world economy.

The lecture sketches the long real price history of oil from the nineteenth century onward. Early price swings were associated with discovery and new uses. The stable middle decades are then explained by a now-vanished institutional arrangement: the United States, especially Texas, was a major oil producer, and the Texas Railroad Commission acted to stabilize prices. There was, in that interval, little need for a highly developed oil futures market.

The world then changed. The first oil crisis of 1973--74 marked a dramatic break. Shiller links it to the Yom Kippur War, to OPEC, and to the broader fact that oil production and pricing power had shifted away from the United States. He also adds a second historical layer: nationalization. Oil that had once been controlled by international oil companies was progressively taken back by sovereign states, beginning early with Mexico in 1938 and later with countries such as Iran. The result was a new geopolitical structure around oil supply, and then a first modern oil crisis.

The second oil crisis of 1979--80, associated with the Iranian Revolution and the Iran--Iraq War, reinforced the lesson. So did later war-related disruptions and finally the great price spike of 2008. Shiller's large claim is that every major oil shock created a recessionary disturbance because the world economy was unprepared for abrupt changes in the price of energy. That is why oil futures are not a minor financial curiosity. They are one of the world's main risk-management institutions.

The lecture then returns, briefly but pointedly, to the logic of safety. If one had locked in an oil price through futures before a major spike, the spike itself would no longer determine one's purchase price. The market is ``civilized,'' as the lecture puts it, not because the world is orderly, but because daily settlement and exchange discipline keep the contract reliable even while geopolitics is violent.

The final move is to financial futures, especially the S\&P 500 index futures market. Here the delivery mechanism changes. We no longer ship wheat or oil to a warehouse. We settle in cash. But the carrying-cost logic survives, altered to fit the nature of the underlying.

Stocks do not generate physical storage expense. Instead they pay dividends. In the lecture's language, storage cost is effectively negative. If we keep the same notation and interpret $y$ now as the dividend yield over the contract horizon, the stock-index analogue of fair value is
\begin{equation}
P_F \approx (1+r-y)P_{\text{spot}}.
\end{equation}
This is again a cautious reconstruction. The transcript is slightly garbled at this point, but the intended structure is clear: dividends offset financing cost.

The lecture then gives the key intuition. If short interest rates are very low and dividend yields are around $2\%$, then short-horizon stock-index futures may naturally lie below spot. At longer horizons the relationship may reverse. But the really important claim is not the sign. It is the interpretation. Because stocks are always ``in storage,'' the S\&P 500 futures curve mostly reflects fair value. It is not a deep forecast of the future of the stock market. Gold is similar in the lecture's closing comparison: it is always stored, so its futures curve often says much more about carrying cost than about coming scarcity.

This is the lecture's final distinction. Some futures markets are highly informative about storage, scarcity, and operational pressure. Others are mainly carrying-cost machines.

\section{Summary}

The chapter begins with the future and ends with carrying cost, but the argument is continuous from start to finish. Futures markets matter because they give us prices for future conditions, and those prices enter the calculations by which modern life is organized.

Agriculture makes the logic visible. Grain must be stored; storage decisions today alter prices and consumption today; and speculation, far from being automatically destructive, can help smooth scarcity over time. The move from forward contracts to futures contracts then shows why standardization, centralized trading, and exchange discipline matter. They solve the problems of heterogeneous terms, weak price discovery, and counterparty risk.

Only after that institutional groundwork does the lecture write the fair-value benchmark
\[
P_F=(1+r+s)P_{\text{spot}},
\]
with $r$ the total interest cost to maturity and $s$ the storage cost over the same horizon. That benchmark explains ordinary contango and the convergence of futures to spot at expiration. The seasonal calendar-time sketch then clarifies how the same logic looks when we think in actual time rather than in maturity space.

Oil supplies the crucial refinement. A downward-sloping futures curve does not overthrow fair value; it tells us that the simple storage relation is not the whole story. Once inventory carries a convenience yield, backwardation can appear. The final extension to stock-index futures then shows a different limit: when the underlying is always effectively in storage, the futures curve may reveal little more than fair value itself.

The lecture therefore leaves us with a coherent picture. Futures markets are not merely speculative side arenas. They are institutions in which history, storage, standardization, and mathematics meet, and they are among the main ways an economy learns to put a price on time.